% ============================================================
%  chapters/appendix_A.tex  — Data Appendix
% ============================================================

\section{Data Appendix}
\label{app:data}

\subsection{Variable Definitions}
\label{app:variable_defs}

Table~\ref{tab:vardefs} defines all variables used in the analysis.

\begin{longtable}{p{0.22\textwidth}p{0.70\textwidth}}
    \caption{Variable Definitions} \label{tab:vardefs} \\
    \toprule
    Variable & Definition \\
    \midrule
    \endfirsthead
    \multicolumn{2}{c}{\tablename~\thetable{} (continued)} \\
    \toprule
    Variable & Definition \\
    \midrule
    \endhead
    \midrule
    \multicolumn{2}{r}{\textit{Continued on next page}} \\
    \endfoot
    \bottomrule
    \endlastfoot

    \textbf{viol\_uter\_p}  & Binary. Experienced violence in LGA 17 to 9 months before birth. \\[0.4em]
    \textbf{viol\_utero}    & Binary. Experienced violence in LGA 0 to 9 months before birth (in utero). \\[0.4em]
    \textbf{viol\_ech}      & Binary. Experienced violence in LGA ages 0–3 (early childhood). \\[0.4em]
    \textbf{viol\_ps}       & Binary. Experienced violence in LGA ages 3–6 (preschool). \\[0.4em]
    \textbf{viol\_prim}     & Binary. Experienced violence in LGA ages 6–12 (primary school). \\[0.4em]
    \textbf{viol\_JS}       & Binary. Experienced violence in LGA ages 12–15 (junior secondary). \\[0.4em]
    \textbf{viol\_SS}       & Binary. Experienced violence in LGA ages 15–18 (senior secondary). \\[0.4em]
    \textbf{attend}         & Binary. Equal to 1 if student is currently enrolled in school. \\[0.4em]
    \textbf{highest\_enrolled} & Continuous. Total years of schooling completed (constructed from school level and class reached). \\[0.4em]
    \textbf{mom\_enrol}     & Continuous. Mother's years of schooling. \\[0.4em]
    \textbf{dad\_enrol}     & Continuous. Father's years of schooling. \\[0.4em]
    \textbf{dist\_JSS}      & Continuous. Distance in kilometers to the junior secondary school the child attends. \\[0.4em]
    \textbf{dist\_SSS}      & Continuous. Distance in kilometers to the senior secondary school the child attends. \\[0.4em]
    \textbf{dist\_prim}     & Continuous. Distance in kilometers to the primary school the child attends. \\[0.4em]
    \textbf{late}           & Binary. Equal to 1 if student began school after reaching age 6. Constructed from survey questions on reasons for late school entry.$^{a}$ \\[0.4em]
    \textbf{age2}           & Continuous. Age-squared. \\

\end{longtable}

\noindent$^{a}$ The dataset contained an age-at-entry variable that appeared
erroneous in some cases (e.g., reported ages younger than the student's
current age, or reported starting age consistent with late entry despite no
late-entry reason being given). The \textit{late} variable was constructed
from a series of survey questions asking the student to indicate why they
started school late. See Section~\ref{sec:methodology} for details.

\subsection{UCDP Data Collection Process}
\label{app:ucdp_process}

The Uppsala Conflict Data Program (UCDP) maintains accuracy and
comprehensiveness through a four-step data collection process
\citep{sundberg2010}:

\begin{enumerate}
    \item An online search is conducted using Factiva, an online database
    carrying more than 10,000 newswires, newspapers, and other sources.
    For each country, at least one major newswire (Reuters, AFP, Xinhua, EFE)
    and BBC Monitoring are searched using conflict-specific terms. Results
    are coded according to UCDP criteria.

    \item Other available materials are searched and coded, including
    recently published books, case studies, journals, NGO publications, and
    other online databases.

    \item Regional experts are consulted to clarify remaining ambiguities.

    \item Finally, UCDP project managers and program directors review the
    data and determine what is included in the final dataset.
\end{enumerate}

\subsection{UCDP Violence Definitions}
\label{app:violence_defs}

\begin{description}
    \item[\textbf{Armed force}] Use of arms to promote a party's general
    position in conflict, resulting in deaths.

    \item[\textbf{Arms}] Any material means, including manufactured weapons
    or sticks, stones, fire, water, etc.

    \item[\textbf{Organized actor}] One of the following: (a) the government
    of an independent state (the party controlling the capital); (b) a
    formally organized group with an announced name using armed force against
    a government, another formally organized group, or civilians; (c) an
    informally organized group without an announced name that uses armed force
    against similarly organized groups in a clear pattern of connected violent
    incidents.

    \item[\textbf{Direct death}] A death directly related to combat between
    warring parties or against civilians.

    \item[\textbf{Civilians}] Unarmed populations that are not part of the
    organized actors.
\end{description}

\paragraph{Event types:}
\begin{itemize}
    \item \textbf{Single-day}: A single incident of armed violence, or
    separate incidents resulting in at least one fatality occurring within
    one 24-hour calendar day at a specified location. Multiple incidents at
    the same location on the same day involving the same actors may be
    aggregated.
    \item \textbf{Summary}: Deaths from multiple separate events over
    several days for which disaggregated information is unavailable.
    \item \textbf{Continuous}: Armed activities at a specified location
    covering multiple interrelated days of continuous fighting.
\end{itemize}

\paragraph{Violence types:}
\begin{itemize}
    \item \textbf{State-based}: Conflicts between two organized groups where
    at least one is the state government.
    \item \textbf{Non-state}: Conflicts between two organized armed groups
    where neither is the state government.
    \item \textbf{One-sided}: Conflicts between an organized actor and
    civilians.
\end{itemize}
